\documentclass[11pt]{article}
\usepackage[utf8]{inputenc}
\usepackage[T1]{fontenc}
\usepackage{lmodern}
\usepackage[margin=1in]{geometry}
\usepackage{amsmath,amssymb,amsthm,mathtools}
\usepackage{booktabs}
\usepackage{enumitem}
\usepackage[colorlinks=true,linkcolor=black,citecolor=black,urlcolor=black]{hyperref}
\usepackage{longtable}
\usepackage{array}
\usepackage{graphicx}
\usepackage{microtype}
\usepackage[most]{tcolorbox}
\tcbuselibrary{skins,breakable}
\usepackage{setspace}
\setstretch{1.05}
\setlength{\parskip}{0.65em}
\setlength{\parindent}{0pt}

\newtheorem{theorem}{Theorem}[section]
\newtheorem{definition}[theorem]{Definition}
\newtheorem{remark}[theorem]{Remark}
\newtheorem{proposition}[theorem]{Proposition}
\newtheorem{lemma}[theorem]{Lemma}

\newcommand{\coderef}[2]{\texttt{#1}\,(\texttt{#2})}

\newtcolorbox{apfbox}{
  breakable, enhanced,
  colback=black!3, colframe=black!30,
  boxrule=0.6pt, arc=2pt,
  left=8pt, right=8pt, top=8pt, bottom=8pt
}

\newtcolorbox{wayfinder}{
  breakable, enhanced,
  colback=blue!3, colframe=blue!40,
  boxrule=0.7pt, arc=2pt,
  left=8pt, right=8pt, top=8pt, bottom=8pt,
  title={\textbf{Wayfinder}},
  fonttitle=\bfseries,
  coltitle=blue!50!black, colbacktitle=blue!10
}

\newtcolorbox{bridgebox}{
  breakable, enhanced,
  colback=orange!3, colframe=orange!50,
  boxrule=0.7pt, arc=2pt,
  left=8pt, right=8pt, top=8pt, bottom=8pt,
  title={\textbf{Bridge}},
  fonttitle=\bfseries,
  coltitle=orange!60!black, colbacktitle=orange!10
}

\title{Paper 5 --- REPRESENTATIONS (Quantum):\\
Quantum Structure as Least-Cost Admissible Bookkeeping\\[2pt]
\large Hilbert Space, the Born Rule, and Measurement under Finite Enforceability}
\author{E. Brooke}
\date{April 2026 (v2.0-PLEC)}

\begin{document}
\maketitle

\begin{abstract}
This paper derives the mathematical structure of quantum mechanics --- Hilbert space,
the Born rule, tensor products, unitarity, and CPTP dynamics --- as the realized
least-cost admissible bookkeeping for coherent jointly enforceable distinctions under
finite capacity. Quantum mechanics is not postulated; it is selected by the
admissibility framework of Papers~1--4 in regimes where coherent reversible refinement
holds. The derivation proceeds without assuming probability, dynamics, or any quantum
postulates. The v2.0-PLEC release reframes the entire derivation around the Principle
of Least Enforcement Cost (PLEC): admissibility fixes the licensable class of
bookkeeping schemes; PLEC selects the realized minimal scheme; measurement is a
change of admissible class rather than a collapse axiom. Imported closure theorems
(Sol\`er, Gleason) are tagged as imports closing an APF-prepared admissible class,
not as APF-internal derivations.
\end{abstract}

\begin{wayfinder}
\textbf{Operational primitives used in this paper.}
\begin{itemize}[leftmargin=1.5em,nosep]
    \item \textbf{Admissible class $\mathcal{A}_{\Gamma}$:} the bookkeeping schemes
    physically licensable in context $\Gamma$ under A1.
    \item \textbf{Enforcement-cost functional $E_{\Gamma}$:} total capacity burden
    required to maintain a configuration $X \in \mathcal{A}_{\Gamma}$.
    \item \textbf{PLEC selection $X^{\ast}$:} realized structure, $X^{\ast} \in
    \arg\min_{X \in \mathcal{A}_{\Gamma}} E_{\Gamma}(X)$.
    \item \textbf{Coherent regime:} reversible refinement, recombination, undo all
    preserve admissibility.
    \item \textbf{Record-locked regime:} capacity irreversibly committed to a logged
    alternative; algebra extends to $M_{\mathrm{sys}} \otimes Z_R$.
\end{itemize}
This paper writes the quantum formalism as the realized PLEC selector on the coherent
admissible class, and writes measurement as the regime exit from coherent to
record-locked admissible bookkeeping.
\end{wayfinder}

\section*{Shared Canonical Frame (Papers 5--6)}

\begin{apfbox}
\textbf{Meaning, Admissibility, Realization, Regime Exit.}
\begin{enumerate}[leftmargin=1.6em]
    \item \textbf{Meaning.} A distinction is physically meaningful iff it is robustly
    enforceable under admissible perturbations.
    \item \textbf{Admissibility.} A configuration, path, or representation is
    admissible iff its total enforcement demand does not exceed available capacity.
    \item \textbf{Realization.} When the admissible class is nonempty and the
    enforcement-cost functional is well-defined on it, realized physical structure
    is selected as the least-enforcement-cost member of that admissible class.
    \item \textbf{Regime exit.} When the conditions required for comparison by
    smooth finite cost fail, the representation exits its domain, while the deeper
    admissibility constraints remain in force.
\end{enumerate}
\end{apfbox}

\section{Introduction}

\subsection{The Question}
Why does quantum mechanics have the mathematical structure it does? Standard
presentations take Hilbert spaces, complex amplitudes, the Born rule, tensor
products, and unitary evolution as axioms or postulates --- accepted because they
work. But the formalism permits vastly more structure than nature realizes. Why the
specific subset that is physically realized?

\subsection{The Answer}
Quantum structure is not arbitrary. It is the least-cost admissible bookkeeping
required to track coherent jointly enforceable distinctions under finite
enforceability. The admissibility framework of Papers~1--4 determines which
bookkeeping structures are physically licensable in coherent regimes; quantum
mechanics is the realized minimal structure within that admissible class. In this
paper we show that, in bilinear unsaturated regimes, least-cost admissible
bookkeeping has the familiar form: linear structure from interference closure,
inner products from admissibility-invariant overlap, complex structure from
continuous reversible symmetry, tensor products from compositional closure, the
Born rule from admissibility-invariant probability assignment, and unitary/CPTP
evolution from admissibility-preserving reversible/general transformations. No
quantum postulates are assumed. The formalism is selected by admissibility.

\subsection{What This Paper Assumes}
This paper assumes Papers~1 (SPINE), 2 (STRUCTURE), 3 (LEDGERS), 4 (CONSTRAINTS).
It also assumes a bilinear unsaturated regime: the interaction functional
$I_{\Gamma}$ is approximately bilinear and no admissibility saturation is
encountered. This is a regime condition, not a universal truth.

\subsection{What This Paper Does NOT Assume}
Hilbert space (derived), complex numbers (derived), Born rule (derived), unitarity
(derived), measurement postulates (regime transitions, not postulates). No quantum
axioms are assumed.

\subsection{Reading Guide}
Section~\ref{sec:plec} introduces the canonical PLEC schema. Sections~\ref{sec:linear}--\ref{sec:dynamics}
walk through the realized least-cost bookkeeping structure step by step.
Section~\ref{sec:measurement} treats measurement as regime exit. Section~\ref{sec:matter}
links the bookkeeping derivation to the matter-content results of Paper~4.
Section~\ref{sec:summary} closes with an explicit epistemic-tag table separating
internally derived results [P], imported closures [P]+[I], and structural
identifications [P\_structural].

\section{Canonical PLEC Schema}\label{sec:plec}

\begin{definition}[Admissible bookkeeping class]
Let $\Gamma$ denote a physical context, interface, or regime. Let $\mathcal{A}_{\Gamma}$
denote the class of all bookkeeping schemes $X$ that satisfy the APF admissibility
conditions in that context.
\end{definition}

\begin{definition}[Enforcement-cost functional]
Let $E_{\Gamma} : \mathcal{A}_{\Gamma} \to [0,\infty)$ assign the total capacity
burden required to maintain $X$ in context $\Gamma$. The minimal value
$\inf_{X \in \mathcal{A}_{\Gamma}} E_{\Gamma}(X)$ defines the realized cost in
$\Gamma$.
\end{definition}

\begin{definition}[PLEC selection]
A realized bookkeeping structure $X^{\ast}$ in context $\Gamma$ is selected by
\[
X^{\ast} \in \operatorname*{arg\,min}_{X \in \mathcal{A}_{\Gamma}} E_{\Gamma}(X).
\]
When the minimizer is unique up to the relevant equivalence relation $\sim$ (basis,
gauge, descriptive coding redundancy), the realized structure is unique up to
$\sim$.
\end{definition}

\begin{remark}
The PLEC schema separates what is \emph{physically licensable} from what is
\emph{physically realized}. The first is fixed by admissibility; the second is
fixed only when the admissible class supports a meaningful least-cost comparison.
The remainder of this paper instantiates the schema for coherent unsaturated
regimes and shows that the realized bookkeeping is quantum mechanics.
\end{remark}

\section{Coherence Defines the Admissible Bookkeeping Class}\label{sec:linear}

\subsection{Coherent Distinctions}
A set of distinctions is \emph{coherent} when it can be reversibly refined,
recombined, and undone without violating admissibility. Coherence does more than
describe a convenient regime: it defines the class of bookkeeping structures that
are physically licensable in unsaturated reversible settings. A bookkeeping scheme
that fails coherence is not merely inconvenient; it does not correctly represent
the admissible refinement structure of the underlying distinctions.

\subsection{Interference Closure}
Interference closure is the central consistency condition on the coherent
admissible class. If two refinement paths lead to the same physical situation,
bookkeeping must identify them consistently. Coherent bookkeeping is therefore
required to close under admissible recombination. Concretely: a distinction $d$
that can be refined two ways --- $d \to \{d_1, d_2\}$ and $d \to \{d_1', d_2'\}$
--- must yield bookkeeping that is closed under reversible recombination of the
two refinements.

\subsection{Linearity}
\begin{theorem}[Interference closure forces linear structure]
Within the coherent admissible class, interference closure forces the space of
realized bookkeeping states to form a vector space over a division algebra.
\end{theorem}
This is the first PLEC selection. Quantum linearity is not postulated in advance;
it is the minimal closure property of admissible coherent bookkeeping. Among all
licensable closure patterns, only the linear closure satisfies coherent
recombination at minimum redundancy.

\section{Invariant Overlap and the Selection of Complex Hilbert Structure}\label{sec:hilbert}

\subsection{Admissibility-Invariant Overlap}
Once coherent bookkeeping is restricted to the admissible class, further structure
is fixed by invariance requirements. A similarity measure between coherent
realizations must remain invariant under admissible reversible transformations and
distribute appropriately over refinements. This forces a sesquilinear form: an
inner product. Completing that structure yields a Hilbert space.

\subsection{Why Complex Numbers?}
Admissibility alone does not license every imaginable Hilbert-like representation.
The admissible class is further restricted by continuity of reversible symmetry.
By Sol\`er's theorem, infinite-dimensional orthomodular spaces over a division
algebra reduce to $\mathbb{R}, \mathbb{C}, \mathbb{H}$. Physical continuity of
one-parameter reversible families in two-dimensional coherent subspaces selects
$\mathbb{C}$. The phase $e^{i\theta}$ parameterizes the unique continuous family
of admissibility-preserving rotations in a 2D coherent subspace. Sol\`er is cited
as an external import; there is no standalone APF check for it, but the
APF-internal preparation of the orthomodular class is the ST-chain in
\coderef{check\_L\_ST\_algebra}{core.py} and its successors (\texttt{check\_L\_ST\_Hilbert},
\texttt{check\_L\_ST\_Dirac}).

\begin{bridgebox}
The role of Sol\`er here is precisely a closure theorem on an APF-prepared
admissible class. APF supplies the orthomodular structure, the division-algebra
structure, and the continuity requirement on one-parameter reversible families;
Sol\`er then closes the classification to $\mathbb{R}, \mathbb{C}, \mathbb{H}$;
continuous symmetry selects $\mathbb{C}$. Tag: [P] + [I: Sol\`er].
\end{bridgebox}

\subsection{Hilbert Space}
Completing the inner-product space yields a Hilbert space $\mathcal{H}$. This is
the realized least-cost admissible representation of coherent distinction
structure in regimes where reversible refinement and continuous symmetry both
hold. The complex Hilbert space is therefore not an arbitrary mathematical
preference; it is the minimum-redundancy bookkeeping consistent with admissible
coherent refinement.

\section{Composition and the Realized Composite Structure}\label{sec:tensor}

\subsection{Joint Enforceability is Conditional}
Papers~1--2 establish that composition can fail: not every pair of systems is
jointly enforceable. This means the admissible class of composite bookkeeping
structures is already conditional. When two systems are jointly enforceable,
bookkeeping must represent that success of composition without violating
refinement consistency.

\subsection{Bilinear Composition Forces the Tensor Product}
Let $A$ and $B$ be coherent systems with Hilbert spaces $\mathcal{H}_A$ and
$\mathcal{H}_B$. If they are jointly enforceable, refinement composability
requires that refining $A$ then composing with $B$ gives the same result as
composing then refining. Bilinearity follows. The minimal admissible composite
bookkeeping satisfying bilinear composition and closure under admissible
recombination is the tensor product $\mathcal{H}_A \otimes \mathcal{H}_B$.

\subsection{Tensor Products are Realized, Not Universal}
Tensor products are not universally primitive. They are the realized least-cost
bookkeeping for composition \emph{within the admissible region}. Outside that
region --- at saturation, measurement, or horizons --- the tensor-product
representation itself ceases to apply. Quantum composition is therefore
conditional at the same level as physical joint enforceability.

\subsection{Entanglement}
Most states in $\mathcal{H}_A \otimes \mathcal{H}_B$ are not product states. These
are entangled states: joint realizations that cannot be factored into independent
realizations. Entanglement is not a mystery in this framework. It is the generic
representation of joint enforceability when the interaction functional is nonzero;
entangled states encode correlations that require coordinated bookkeeping across
interfaces.

\subsection{Conditional Composition (Paper 2 Link)}
The tensor product exists \emph{if and only if} the composite is admissible:
$A \otimes B \text{ exists} \iff E_{\Gamma}(A \cup B) \le C_{\Gamma}$. This is
non-closure at the level of quantum representations: the familiar assumption that
composition always works is valid only within the admissible region. At saturation
boundaries, the tensor-product description itself breaks down and the regime
exits.

\section{The Born Rule as the Unique Admissible Probability Assignment}\label{sec:born}

\subsection{Refinement Invariance}
Probability assignments must be refinement-invariant: coarse-graining and
fine-graining must agree. If outcome $E$ can be refined into $E_1 \cup E_2$ with
$E_1 \cap E_2 = \emptyset$, then $p(E) = p(E_1) + p(E_2)$. This is additivity over
disjoint outcomes.

\subsection{Admissibility Invariance}
Probability assignments must be admissibility-invariant: they must not depend on
arbitrary bookkeeping choices. For a state $\rho$ and effect $E$, $p(E\mid\rho)$
must be invariant under basis changes, equivalent representations, and admissible
reversible transformations.

\subsection{Gleason's Theorem}
\begin{theorem}[Gleason 1957, as a closure on the admissible class]
In a complex Hilbert space of dimension $\geq 3$, the unique probability measure
satisfying refinement additivity and admissibility-invariance is
\[
p(E\mid\rho) = \mathrm{Tr}(\rho E).
\]
For pure states $\rho = \lvert\psi\rangle\langle\psi\rvert$ and projections
$E = \lvert\phi\rangle\langle\phi\rvert$, this gives $p = \lvert\langle\psi\mid\phi\rangle\rvert^2$.
\end{theorem}
APF-internal verification of the finite-dimensional case is in
\coderef{check\_L\_Gleason\_finite}{supplements.py}; the full Gleason statement
remains an external import.

\begin{bridgebox}
The important point is not merely that Gleason's theorem applies, but \emph{why}
it applies here. APF does not conjure probability from nowhere; it identifies the
admissible class of bookkeeping structures and the invariance conditions any
physically meaningful assignment must satisfy. The Born rule is then the unique
probability assignment compatible with that realized bookkeeping. Tag: [P] + [I:
Gleason 1957].
\end{bridgebox}

\subsection{The Born Rule Is Selected, Not Postulated}
The Born rule is not a postulate. It is the unique admissibility-invariant
probability assignment in the coherent admissible class once Hilbert structure is
selected. The ``mystery'' of why probability is $\lvert\psi\rvert^2$ dissolves: it
is the only assignment compatible with the bookkeeping forced by admissibility.

\section{Dynamics as Admissibility-Preserving Evolution}\label{sec:dynamics}

\subsection{Reversible Evolution: Unitarity}
\begin{theorem}[Admissibility-preserving reversible evolution is unitary]
Within the realized coherent bookkeeping structure, reversible evolution must
preserve admissibility. Preservation of the inner product forces unitary or
antiunitary transformations; physical continuity excludes antiunitary maps except
for discrete cases such as time reversal. Continuous reversible evolution is
therefore generated by self-adjoint operators: $U(t) = e^{-iHt}$. This is the
Schr\"odinger structure.
\end{theorem}

\subsection{General Evolution: CPTP Maps}
Not all evolution is reversible. Record formation, environmental interaction, and
capacity commitment produce irreversible changes. Admissibility-preserving
transformations that remain admissibility-preserving under extension by ancillas
are completely positive. With normalization (trace preservation),
admissibility-preserving evolution is described by CPTP maps.

\subsection{Dynamics is Not a Separate Postulate}
The familiar dynamical structure of quantum mechanics is not a separate postulate
layered on top of Hilbert space. It is the admissibility-preserving evolution law
of the realized coherent bookkeeping selected in the preceding sections.

\section{Measurement as Exit from the Coherent Bookkeeping Regime}\label{sec:measurement}

\subsection{The Measurement Problem Re-Read}
Measurement is the central conceptual difficulty of standard quantum mechanics
because the formalism appears to require an extra collapse postulate that
interrupts otherwise unitary evolution. In the admissibility framework, the issue
is reformulated more cleanly. The question is not how a unitary system suddenly
violates its own law, but \emph{when the coherent bookkeeping regime ceases to be
the correct admissible representation}. Measurement is therefore a
regime-transition problem, not an axiom problem.

\subsection{Before and After Measurement as Class Change}
Before measurement, multiple outcomes are jointly admissible and the relevant
bookkeeping class is the coherent class: reversible refinement, recombination, and
admissibility-invariant overlap are all meaningful. After measurement, capacity
has been irreversibly committed to a recorded alternative. The relevant
bookkeeping class is no longer the coherent one; it is the record-locked class in
which mutually incompatible alternatives are excluded by the structure of the
irreversible log.

The key point is that this is a \emph{change of admissible class}. Quantum
superposition does not ``mysteriously disappear.'' Rather, the system exits the
coherent least-cost bookkeeping regime and enters a record-locked regime whose
admissible structure is different.

\subsection{Record-Locking Channels as Explicit Algebraic Class Change (T9 / L3-$\mu$)}
The record-locking channel formalization makes this change explicit. The
interface algebra decomposes as
\[
M = M_{\mathrm{sys}} \otimes Z_R,
\]
where $M_{\mathrm{sys}}$ is the coherent system algebra and $Z_R$ is a classical
log algebra with projections $\{p_s\}$ indexed by log strings. Each enforcement
operation $\mathcal{E}_i$ acts both on the system and through an irreversible
append map $\alpha_i: p_s \mapsto p_{(s,i)}$ on the log:
\[
\mathcal{E}_i = (\mathcal{E}_i^{\mathrm{sys}} \otimes \mathrm{id}) \circ
(\mathrm{id} \otimes \alpha_i).
\]

\begin{theorem}[T9 / L3-$\mu$, $k!$ inequivalent histories]
For $k$ enforcement operations applied in all $k!$ orderings, the resulting CP
maps are pairwise distinct: for $\pi \neq \sigma$, $\mathcal{E}_{\pi} \neq
\mathcal{E}_{\sigma}$. The log projections $p_{\pi}$ and $p_{\sigma}$ are
orthogonal. Therefore $\lvert H_k / \equiv\rvert = k!$ inequivalent histories.
\end{theorem}
Formal verification: \coderef{check\_T9}{gauge.py}.

This construction does more than reinterpret measurement. It exhibits the new
admissible class algebraically. What is traditionally called collapse is the
point at which the physically relevant admissible description is no longer
supported on $M_{\mathrm{sys}}$ alone, but on $M_{\mathrm{sys}} \otimes Z_R$ with
irreversible sector separation. The log is irreversible because record-lock
capacity constraints prevent undoing the append (T0.4$'$ certificate, Paper~2).

\subsection{Decoherence as Progressive Class Transfer}
Decoherence is the progressive transfer from coherent admissible bookkeeping to
record-locked admissible bookkeeping. As environmental records accumulate, the log
algebra $Z_R$ grows and the coherent sector becomes a less accurate effective
compression of the admissible structure. When enough capacity has been committed
to records, the quantum description ceases to apply --- not because it is false,
but because the system is no longer in the regime whose admissible class quantum
mechanics represents.

\begin{bridgebox}
Measurement is to Paper~5 what saturation, branching, and horizons are to
Paper~6: not a violation of deeper law, but a boundary of representational
validity. The unifying project-level statement is: record locking, saturation,
horizons, and topology change are all regime exits of representation, not
failures of admissibility itself.
\end{bridgebox}

\section{Connection to Matter Content}\label{sec:matter}

\subsection{Quantum Structure Constrains Particle Content}
The quantum representation derived in this paper does not float freely above the
constraint structure. It is jointly constrained with the gauge and matter content
derived in Papers~1 and 4 by the same capacity architecture:

\begin{center}
\begin{tabular}{@{}llll@{}}
\toprule
Quantum structure & Capacity constraint & Consequence & Coderef \\
\midrule
$\mathcal{H}_A \otimes \mathcal{H}_B$ & Conditional on admissibility & Composition fails at horizons & --- \\
Unitary symmetry group & Gauge selection & $SU(3)\!\times\!SU(2)\!\times\!U(1)$ & \coderef{check\_T\_gauge}{gauge.py} \\
Matter representations & Anomaly cancellation & Chiral fermion content forced & \coderef{check\_T5}{gauge.py} \\
Generation structure & Capacity saturation & $N_{\mathrm{gen}} = 3$ & \coderef{check\_T4F}{gauge.py} \\
Born-rule probabilities & Gleason + admissibility & $p = \lvert\langle\psi\mid\phi\rangle\rvert^2$ & \coderef{check\_L\_Gleason\_finite}{supplements.py} \\
\bottomrule
\end{tabular}
\end{center}

The quantum representation is not free to have arbitrary matter content. The same
capacity architecture that forces Hilbert space also forces three generations of
chiral fermions in the Standard Model gauge group. This is qualitatively new:
prior derivations of quantum structure (Hardy, Chiribella et al.) do not constrain
particle content.

\subsection{Particles Require Quantum Structure}
The enforcement potential $V(\Phi)$ derived in Paper~3 (Section~5.3) shows that
particle-like excitations require the quantum structure derived here. The binding
well at $\Phi/C \approx 0.81$ has mass gap $d^2 V/d\Phi^2 > 0$, but no classical
soliton localizes: particles are quantum modes of the enforcement field, not
classical cost minima. This closes the conceptual gap between the constraint
architecture and observable particle physics.

\subsection{Bridge Tightening (Open)}
The bridge from bookkeeping structure to gauge and matter content is presently a
structural coexistence claim under common capacity architecture, not a direct
single-theorem derivation. Tightening the bridge from compatibility to realized
selection is one of the open theorem projects identified in the v2.0-PLEC
revision program.

\section{Epistemic Tagging Discipline (Paper 5)}\label{sec:tagging}

\begin{apfbox}
\textbf{Project-wide categories.}
\begin{itemize}[leftmargin=1.5em,nosep]
    \item \textbf{[P] Internally derived.} Result follows from APF axioms plus
    explicitly stated regime assumptions, without external classification.
    \item \textbf{[P]+[I] Derived with import.} APF fixes the admissible class; an
    external mathematical theorem closes classification, uniqueness, or representation.
    \item \textbf{[P\_structural] Structural identification.} The framework
    identifies an already-derived structure with a familiar physical role; not a
    stronger derivation than the identification supports.
    \item \textbf{[Pred] Phenomenological proposal/prediction.} Empirical
    identification or scaling whose status is downstream and must be evaluated
    against data.
\end{itemize}
\end{apfbox}

\begin{center}
\begin{tabular}{@{}llll@{}}
\toprule
Item & Tag & Coderef & Reason \\
\midrule
Linear structure & [P] & --- & Interference closure on the coherent admissible class. \\
Inner product & [P] & --- & Admissibility-invariant overlap. \\
Complex field selection & [P]+[I] & \coderef{check\_L\_ST\_Hilbert}{core.py} & APF narrows; Sol\`er closes classification. \\
Hilbert completion & [P] & \coderef{check\_L\_ST\_Hilbert}{core.py} & Completion of inner-product structure. \\
Tensor product & [P] & --- & Minimal admissible composite under bilinearity. \\
Born rule & [P]+[I] & \coderef{check\_L\_Gleason\_finite}{supplements.py} & APF fixes invariance; Gleason closes. \\
Unitarity & [P] & --- & Inner-product preservation + continuity. \\
CPTP maps & [P] & --- & Admissibility preservation under ancilla extension. \\
Record-locking, $k!$ histories & [P] & \coderef{check\_T9}{gauge.py} & Explicit algebraic construction. \\
Gauge coupling & [P] & \coderef{check\_T\_gauge}{gauge.py} & $SU(3)\!\times\!SU(2)\!\times\!U(1)$ selection. \\
Generation count & [P] & \coderef{check\_T4F}{gauge.py} & $N_{\mathrm{gen}} = 3$ from saturation. \\
Anomaly cancellation & [P] & \coderef{check\_T5}{gauge.py} & Chiral fermion content forced. \\
Measurement as class change & [P\_structural] & \coderef{check\_Regime\_exit\_Type\_III}{plec.py}, \coderef{check\_T9}{gauge.py} & Regime transition (v6.9 formalization), not primitive axiom. \\
Matter-content bridge & mixed & multiple & Structural coexistence; tightening open. \\
\bottomrule
\end{tabular}
\end{center}

This discipline makes the paper easier to defend: imported steps remain visible,
while internally derived structure becomes sharper rather than weaker.

\section{Summary}\label{sec:summary}

\subsection{What Is Derived}
The realized least-cost admissible bookkeeping for coherent jointly enforceable
distinctions in bilinear unsaturated regimes:
\begin{itemize}[leftmargin=1.5em,nosep]
    \item Linear vector-space structure from interference closure.
    \item Inner-product structure from admissibility-invariant overlap.
    \item Complex Hilbert space via Sol\`er + continuous symmetry selection.
    \item Tensor-product composition (conditional on joint admissibility).
    \item The Born rule via Gleason on the prepared admissible class.
    \item Unitary and CPTP evolution as admissibility preservation.
    \item Record-locking and the $k!$ history sectors of T9 / L3-$\mu$.
\end{itemize}

\subsection{What Is NOT Derived}
Specific Hamiltonians (depend on physical content). Coupling constants beyond
the structural ratios (regime parameters). Particle masses beyond the structural
hierarchy (Yukawa ratios are regime parameters). Particle identification of dark
matter (Paper~6 fences this).

\subsection{The Status of Quantum Mechanics}
Quantum mechanics is the realized least-cost bookkeeping for coherent jointly
enforceable distinctions in bilinear unsaturated regimes. It is exact within that
regime, conditional on its regime assumptions, and non-fundamental in the sense
that it is a representation rather than the ontological base layer. When the
conditions of coherent reversible bookkeeping fail --- at saturation, record
locking, measurement, or horizons --- the quantum representation exits its
domain, while the deeper admissibility constraints remain in force.

\appendix

\section*{Appendix A: Regime Assumptions and Where They Fail}

\begin{center}
\begin{tabular}{@{}lll@{}}
\toprule
Assumption & Content & Where it fails \\
\midrule
Interference closure & Refine / recombine / undo preserve admissibility & Record-locked histories; horizons \\
Overlap invariance & Similarity invariant under reversible transforms & Near saturation boundaries \\
Phase admissibility & Continuous one-parameter reversible families exist in 2D subspaces & Discrete / quantized regimes \\
Compositional closure & Refinement and composition commute for jointly enforceable systems & Saturation; horizon crossing \\
Bilinearity & $I_{\Gamma}$ approximately bilinear & Strong-coupling; near-saturation \\
\bottomrule
\end{tabular}
\end{center}
Where these regime assumptions fail, quantum mechanics ceases to apply --- but the
underlying admissibility constraints (Papers~1--4) continue to hold. The
constraint structure is more fundamental than the quantum representation.

\section*{Appendix B: Record-Locking Channel Details}

The T9 / L3-$\mu$ proof constructs the record-locking model explicitly:

\begin{enumerate}[leftmargin=1.6em]
    \item \textbf{Interface algebra.} $M = M_{\mathrm{sys}} \otimes Z_{\le k}$,
    where $Z_{\le k}$ is the classical log algebra with projections $\{p_s\}$ for
    $s \in \Sigma^m$ ($m \le k$), $\Sigma = \{1,\ldots,k\}$.
    \item \textbf{Append maps.} $\alpha_i(p_s) = p_{(s,i)}$ for $\lvert s\rvert <
    k$; $\alpha_i(p_s) = p_s$ for $\lvert s\rvert = k$ (saturation). Key property:
    $\alpha_i(p_s) \cdot \alpha_j(p_s) = 0$ for $i \neq j$ (ranges disjoint).
    \item \textbf{Enforcement channels.} $\mathcal{E}_i = (\mathcal{E}_i^{\mathrm{sys}}
    \otimes \mathrm{id}) \circ (\mathrm{id} \otimes \alpha_i)$. Applying
    $\mathcal{E}_i$ refines the system and appends ``$i$'' to the irreversible log.
    \item \textbf{Permutations.} For permutation $\pi$,
    $\mathcal{E}_{\pi(1)} \circ \cdots \circ \mathcal{E}_{\pi(k)}$ maps
    $1_{\mathrm{sys}} \otimes p_{\emptyset}$ to
    $1_{\mathrm{sys}} \otimes p_{(\pi(1),\ldots,\pi(k))} = 1_{\mathrm{sys}} \otimes p_{\pi}$.
    \item \textbf{Distinctness.} For $\pi \neq \sigma$, $p_{\pi} \neq p_{\sigma}$
    and $p_{\pi} \cdot p_{\sigma} = 0$. Therefore $\lvert H_k / \equiv \rvert = k!$.
\end{enumerate}

\textbf{Corollary.} $\dim(Z_R) \ge k!$ and $\Delta S_{\mathrm{records}} \ge \log(k!)
\sim k\log k$. The entropy of measurement is structural, not thermodynamic.

\section*{Appendix C: Changelog}

\begin{itemize}[leftmargin=1.5em,nosep]
\item \textbf{v2.0-PLEC (April 2026).} Reframed entire derivation around PLEC schema:
admissibility fixes the licensable bookkeeping class, PLEC selects the realized minimum,
measurement is regime exit from coherent to record-locked admissible class. Added shared
canonical box, wayfinder box, canonical PLEC schema (definitions), epistemic tagging
discipline section. Reformatted bridge boxes around Sol\`er and Gleason as imports
closing an APF-prepared admissible class. Cleaned formatting throughout (v1.0 had
literal-escape artifacts in math mode and broken section spacing).
\item \textbf{v1.0 (February 2026).} Record-locking channels (T9 / L3-$\mu$).
Decoherence as progressive capacity commitment. Matter-content connection
(quantum structure jointly constrains particle content). Regime-assumptions
appendix. Forward references to Paper~7 (ACTION) and Paper~4 \S 7
($\sin^2\theta_W$).
\end{itemize}

\end{document}
